We have demonstrated that PHATE, a manifold learning technique originally developed for biological data, generalizes effectively to high-dimensional language embeddings. By leveraging both local and global data geometry, PHATE captures semantic and narrative structures in text with greater clarity than conventional techniques such as PCA, UMAP, and t-SNE. A multiscale variant of PHATE further recovers meaningful hierarchical relationships, outperforming standard pipelines in clustering tasks.

Our results show that the structural irregularities common in real-world text—such as dropout analogs, batch effects, and latent hierarchies—are well-handled by PHATE’s diffusion-based approach. Through both synthetic datasets and real policy narratives, we find that PHATE preserves fine- and coarse-grained semantic structure while enhancing interpretability. These findings highlight PHATE’s potential as a valuable tool for unsupervised text analysis, opening new opportunities for its use in natural language processing and computational social science.